% TikZ source for Figure 1, Panel C (R2 rework v2): a clean 3D coastal-upwelling / process
% schematic. Included via \input from report/sections/02_data_and_case_studies.tex inside a
% \resizebox. The ocean is a connected box (top sea-surface face + south cross-section + west end
% face) so the perspective reads correctly; the coast/headland sits east (right). Process chain:
% alongshore equatorward wind -> offshore surface (Ekman) transport -> upwelling of cold,
% nutrient-rich water -> phytoplankton bloom (Chl a) -> cooler / more variable oxygen, at the
% Tongoy buoy (BTG) / PLV setting hosting scallop aquaculture. Over-water labels carry a white
% halo for legibility; leaders do not cross text. Editable vector source; not to scale.
\definecolor{seaTop}{HTML}{a9cbe0}
\definecolor{seaMid}{HTML}{5b93bd}
\definecolor{seaDeep}{HTML}{1f4f78}
\definecolor{seaEnd}{HTML}{3d78a6}
\definecolor{landC}{HTML}{d8caa4}
\definecolor{landTop}{HTML}{e6dcc0}
\definecolor{landEdge}{HTML}{9c8f66}
\definecolor{upC}{HTML}{1f6fa5}
\definecolor{windC}{HTML}{2B2B2B}
\definecolor{bioC}{HTML}{126b50}
\definecolor{inkC}{HTML}{2B2B2B}
\definecolor{oxyC}{HTML}{7a5e12}

\begin{tikzpicture}[x=0.62cm, y=0.62cm, font=\footnotesize,
    lead/.style={draw=inkC, line width=0.35pt},
    lab/.style={text=inkC, font=\footnotesize, inner sep=1pt},
    wlab/.style={text=inkC, font=\footnotesize, inner sep=1.6pt,
                 fill=white, fill opacity=0.72, text opacity=1, rounded corners=1pt}]

  \def\dx{1.9}\def\dy{1.25}   % perspective depth (into page = alongshore, toward the north)
  \def\S{6.4}                 % sea-surface height (front edge)
  \def\Wc{9.4}                % coast x-position (shoreline)

  % ---- ocean box: west end face, front cross-section, top surface ----
  \fill[seaEnd] (0,\S) -- (\dx,\S+\dy) -- (\dx,2.25) -- (0,1.0) -- cycle;
  \shade[top color=seaMid, bottom color=seaDeep] (0,\S) -- (\Wc,\S) -- (\Wc,3.0) -- (0,1.0) -- cycle;
  \draw[draw=landEdge, line width=0.5pt] (0,1.0) -- (\Wc,3.0);
  \shade[left color=seaTop, right color=seaMid] (0,\S) -- (\Wc,\S) -- (\Wc+\dx,\S+\dy) -- (\dx,\S+\dy) -- cycle;
  \draw[draw=seaDeep, line width=0.3pt, opacity=0.5] (0,\S) -- (\Wc,\S) -- (\Wc+\dx,\S+\dy) -- (\dx,\S+\dy) -- cycle;

  % ---- coastal land: front wedge + headland, thin top strip for 3D ----
  \fill[landC, draw=landEdge, line width=0.5pt]
    (\Wc,\S) -- (10.6,7.1) -- (11.5,8.4) -- (12.7,7.3) -- (15.0,6.9)
    -- (15.0,3.0) -- (\Wc,3.0) -- cycle;
  \fill[landTop, draw=landEdge, line width=0.4pt, opacity=0.9]
    (\Wc,\S) -- (\Wc+\dx,\S+\dy) -- (12.5,8.55) -- (11.5,8.4) -- (10.6,7.1) -- cycle;
  \draw[landEdge, line width=0.6pt] (\Wc,\S) -- (\Wc,3.0);

  % ================= PROCESS ARROWS =================
  % alongshore, upwelling-favourable wind: bold arrows high on the top face (into the page)
  \foreach \sx in {3.3,4.7,6.1}{
    \draw[windC, line width=1.6pt, -{Latex[length=2.6mm]}] (\sx,\S+0.62) -- (\sx+1.35,\S+0.62+\dy*0.72);
  }
  \node[lab, anchor=south, align=center] at (5.6,8.5)
    {\textbf{Equatorward,}\\\textbf{upwelling-favourable wind}};

  % net surface (Ekman) transport -- offshore, at the front edge of the surface
  \draw[upC, line width=1.7pt, -{Latex[length=2.8mm]}] (8.4,\S+0.1) -- (3.2,\S+0.1);
  \node[wlab, anchor=west, align=left, text=upC] at (0.3,\S+0.55) {offshore\\Ekman transport};

  % upwelling of cold, nutrient-rich water near the coast + deep feeder
  \draw[upC, line width=2.0pt, -{Latex[length=3mm]}]
    (8.9,3.3) .. controls (8.4,4.3) and (8.7,5.3) .. (8.5,6.25);
  \draw[upC, line width=1.0pt, opacity=0.7, -{Latex[length=2mm]}] (4.7,1.85) -- (8.4,3.15);
  \node[wlab, anchor=east, align=right, text=upC] at (7.4,2.0) {upwelling of cold,\\nutrient-rich water};

  % phytoplankton bloom (Chl a) along the surface at the coast
  \fill[bioC, opacity=0.34] (6.1,\S-0.4) rectangle (9.25,\S);
  \foreach \bx/\by in {6.5/6.18, 7.1/6.08, 7.7/6.22, 8.3/6.1, 8.85/6.2}{ \fill[bioC] (\bx,\by) circle (1.0pt); }
  \node[wlab, anchor=east, text=bioC, align=right] at (5.95,5.55) {\textbf{phytoplankton}\\\textbf{bloom (Chl\,\textit{a})}};
  \draw[lead, draw=bioC] (5.98,5.72) -- (6.55,6.08);

  % oxygen consequence, cautiously worded, leader to the upwelled column
  \node[wlab, anchor=west, text=oxyC, align=left] at (0.5,3.6) {upwelled water: cooler,\\lower / more variable O\textsubscript{2}};
  \draw[lead, draw=oxyC] (4.0,3.75) -- (7.95,3.95);

  % ================= SITE FEATURES =================
  % Tongoy buoy (BTG)
  \fill[inkC] (8.94,\S) rectangle (9.06,\S+0.72);
  \fill[inkC] (8.8,\S) -- (9.2,\S) -- (9.0,\S-0.22) -- cycle;
  \node[lab, anchor=south west, align=left] at (10.05,7.75) {Tongoy buoy\\(BTG)};
  \draw[lead] (10.05,7.9) -- (9.06,6.7);

  % scallop aquaculture lines just west of the buoy
  \foreach \ax in {7.15,7.45,7.75}{ \draw[inkC, line width=0.7pt] (\ax,\S-0.5) -- (\ax,5.15); \fill[inkC] (\ax,5.1) circle (0.8pt); }
  \node[wlab, anchor=north, align=center] at (7.45,5.0) {scallop\\aquaculture};

  % PLV meteorological station on the headland
  \draw[inkC, line width=0.9pt] (11.5,8.4) -- (11.5,9.25);
  \draw[inkC, line width=0.5pt] (11.3,9.25) -- (11.7,9.25);
  \fill[inkC] (11.5,8.4) circle (1.1pt);
  \node[lab, anchor=west] at (11.75,9.1) {PLV met.\ station};

  % satellite / reanalysis coarse external context (top-left corner)
  \foreach \gx in {0.4,1.0,1.6}{ \draw[inkC, line width=0.3pt, opacity=0.5] (\gx,8.9) -- (\gx,8.25); }
  \foreach \gy in {8.25,8.575,8.9}{ \draw[inkC, line width=0.3pt, opacity=0.5] (0.4,\gy) -- (2.2,\gy); }
  \draw[inkC, line width=0.5pt, opacity=0.6, -{Latex[length=1.8mm]}] (1.3,8.2) -- (1.3,7.45);
  \node[lab, anchor=south west, align=left] at (0.4,8.95) {satellite /\\reanalysis};

\end{tikzpicture}
